\scanpage{191}
\mspage{1}
$S$ schéma, \inserted{$Et(S)$ le site étale de $S$. $\mathcal U$ un univers.}
Une ``donnée de ramification'' sur $S$ est un couple $(\mathscr R,r)$ formé d'un \uline{champ en \inserted{$\mathcal U$-}catégories multigaloisiennes} $\mathscr R$ sur $Et(S)$,
\margin{on sous-entend que les \unsure{fonct$^{\rm s}$|foncteurs} chgt de base sont \uuline{exacts} !}
et d'un morphisme \deleted{de champs} \inserted{cartésien} (``réalisation géométrique'')
\[
r:\mathscr R\longrightarrow \mathrm{Fl}_S
\]
où $\mathrm{Fl}_S$ est la \deleted{\unsure{champ}} catégorie fibrée sur $Et(S)$ dont la fibre en l'objet $S'$ est $(Sch)_{/S'}$ ;
% "catégorie fibrée" is struck through and dotted underneath (stet).
\deleted{\unsure{on donne}} un objet $X$ de $\mathscr R_{S'}$ ( \deleted{\unsure{Et|objet}} $S'\in\mathrm{Ob}\,Et(S))$ sera dit \uline{revêtement de type $\mathscr R$} de $S'$, \deleted{\unsure{et}} $r(X)$ est sa \uline{réalisation} géométrique.
\deleted{On fait l'hypothèse} Ces données sont soumises aux axiomes suivants :

R1 Le foncteur $r$ commute aux sommes finies et au passage au quotient
\inserted{il suffit passage au quotient par groupe fini \unsure{ordinaire} (et \unsure{même|dans les} cas galoisien \unsure{ss pt fixe !})}
par une relation d'équivalence
[il sera $=$ exact à droite, cf \uline{plus bas} \unsure{\ldots}, \uline{mais en général pas exact à gauche} !].
\margin{On veut que dans $\mathrm{Fl}_S$ on ait quotient \uuline{admissible}}

R2 \unsure{Pour tt|\deleted{Donné} Soient} $S'\in\mathrm{Ob}\,Et(S)$ et $X\in\mathrm{Ob}\,\mathscr R_{S'}$, l'application $X'\mapsto r(X')$ est une bijection de l'ens des \deleted{sous-objets}\inserted{sommandes} directs \inserted{(= ss-objets)} $X'$ de $X$ avec l'ens des sommandes directs (i.e. parties ouvertes et fermées) de $r(X)$.

R3 Soit $S'$, $X$ comme dessus, et soit \deleted{\unsure{$(X_i)$}}\inserted{$(X_i)$} une famille (pas néc. finie) de \deleted{\unsure{flèches de but $X$}} sous-objets de $X$, telle que la famille des $r(X_i)$ recouvre $r(X)$. Alors la \deleted{famille $X_i$ admet $X$ comme somme directe effective,}\deleted{\unsure{soit la}} famille $X_i\longrightarrow X$ est une famille de descente effective (= est épimorphique effective \ldots).
% "sous-objets de X" is struck and dotted underneath.

R4 Soit $X\to Y$ dans $\mathscr R_{S'}$ ($S'\in\mathrm{Ob}\,Et(S)$), d'où $Y=\coprod_{n\in\mathbb N}Y_n$, avec $X\times_Y Y_n\to Y_n$ de rang $n$. On veut que les $r(Y_n)$ \uline{recouvrent} $r(Y)$.
\margin{Dans R4, il suffirait de supposer $Y=e$}

\uline{Prop}. \deleted{\unsure{Cor$^{\rm e}$|Écri}} Disons qu'une famille de flèches \inserted{$X_i\longrightarrow X$} de $\mathscr R_{S'}$ est couvrante si la famille $r(X_i)\longrightarrow r(X)$ est surjective, i.e. (grâce : R1 c'est kif kif) si les images $Y_i\subset X$ sont telles que les $r(Y_i)$ recouvrent $r(X)$. Alors ce sont les familles couvrantes d'une \uline{topologie} sur $\mathscr R_{S'}$, \inserted{moins fine que la canonique} \inserted{et pour laquelle tout objet de $\mathscr R_{S'}$ est loc. constant fini.} le foncteur chgt de base $\mathscr R_{S'}\longrightarrow\mathscr R_{S''}$ est continu. \unsure{(donc|Comme|[donc}
\margin{n'utilise que R1 et R3, R4}
% archive sheet number: 190
